% chapters/16/section-05.tex
\section{Log Bug dan Perbaikan}
\label{sec:log-bug}

\subsection{Catatan yang Dapat Diulang}

Sebuah temuan baru berguna apabila orang lain dapat mengulanginya kembali.
``Bunyinya kadang tidak terdengar'' tidak dapat diperbaiki, sebab Anda tidak
tahu kapan ia terjadi. ``Bunyi lompat tidak terdengar apabila pemain melompat
kurang dari setengah detik sesudah mendarat'' dapat diperbaiki pada hari yang
sama.

Perbedaan keduanya terletak pada tiga hal yang harus ada pada setiap catatan
\istilah{bug}: apa yang terjadi, apa yang seharusnya terjadi, dan urutan
tindakan yang membuatnya terjadi. Ketiganya diperlihatkan pada
Tabel~\ref{tab:log-bug}.

\begin{table}[htbp]
\centering
\caption{Log \istilah{bug} beserta kolom yang harus diisi}
\label{tab:log-bug}
{\setlength{\tabcolsep}{4pt}\footnotesize%
\begin{tabularx}{\textwidth}{@{}C{0.8cm} L{2.4cm} Y C{1.8cm} C{1.4cm}@{}}
\toprule
\headrow \textbf{No} & \textbf{Jenis} & \textbf{Urutan yang menyebabkannya, dan hasil yang salah} & \textbf{Berat} & \textbf{Status} \\
\midrule
1 & \istilah{Soft-lock} &
Pemain berjalan ke kiri di ujung lorong, lalu jatuh ke lubang. Seluruh tombol
tidak lagi mengeluarkan pemain dari lubang itu; permainan harus dimulai ulang.
& Berat & Belum \\
\addlinespace
2 & Integrasi &
Bunyi pengambilan benda tidak terdengar pada benda ketiga. Benda pertama dan
kedua berbunyi. Skornya tetap bertambah pada ketiganya.
& Sedang & Belum \\
\addlinespace
3 & Tampilan &
Animasi diam tetap berjalan ketika pemain berlari, apabila berlari dimulai
ketika pemain baru saja mendarat.
& Ringan & Belum \\
\bottomrule
\end{tabularx}}
\end{table}

\subsection{Menimbang Beratnya}

Tidak semua temuan perlu diperbaiki, dan waktu yang tersisa menjelang
pengumpulan selalu terbatas. Karena itu, setiap temuan sebaiknya ditimbang
beratnya, dan penimbangan itu dilakukan atas dua pertanyaan.

\begin{enumerate}
  \item \textbf{Apakah permainannya masih dapat diselesaikan?} Temuan yang
        menjawab ``tidak'' selalu paling berat, walaupun ia jarang terjadi.
        Butir pertama pada Tabel~\ref{tab:log-bug} adalah contohnya.
  \item \textbf{Berapa banyak pemain yang akan menemukannya?} Temuan yang
        dialami setiap pemain pada menit pertama mendesak; temuan yang hanya
        muncul pada keadaan yang jarang dapat ditunda.
\end{enumerate}

\begin{table}[htbp]
\centering
\caption{Skala berat temuan yang dipakai pada praktikum bab ini}
\label{tab:berat-temuan}
{\setlength{\tabcolsep}{4pt}\footnotesize%
\begin{tabularx}{\textwidth}{@{}L{1.8cm} L{4.6cm} Y@{}}
\toprule
\headrow \textbf{Berat} & \textbf{Tandanya} & \textbf{Tindakan} \\
\midrule
Berat & Permainan tidak dapat diselesaikan, atau kemajuannya hilang. &
Diperbaiki sebelum sesi \istilah{playtest} berikutnya. \\
\addlinespace
Sedang & Permainannya selesai, tetapi satu bagian yang dijanjikan GDD tidak
bekerja. &
Diperbaiki, dan diperiksa kembali pada putaran berikutnya. \\
\addlinespace
Ringan & Tampilan atau bunyi yang tidak sesuai, tanpa mengubah jalannya
permainan. &
Diperbaiki apabila waktu tersisa; dicatat sebagai kekurangan yang diketahui. \\
\bottomrule
\end{tabularx}}
\end{table}

\begin{tipsbox}
Kebiasaan yang menghemat waktu: perbaiki temuan berat satu per satu, dan
jalankan \istilah{checklist} fungsional sesudah setiap perbaikan. Menggabungkan
beberapa perbaikan lalu menjalankan pemeriksaan sekali saja membuat Anda tidak
tahu perbaikan mana yang bekerja, dan perbaikan mana yang justru menimbulkan
temuan baru.
\end{tipsbox}

\subsection{Dari Temuan kembali ke GDD}

Lingkaran pengujian belum selesai pada perbaikan. Temuan yang berulang pada
beberapa sesi \istilah{playtest} biasanya bukan kesalahan pelaksanaan,
melainkan tanda bahwa rancangannya memang belum jelas. Apabila setiap penguji
melewati pintu yang belum terbuka, kemungkinan besar GDD Anda tidak pernah
menjelaskan bahwa pintu itu terkunci, sehingga Anda sendiri tidak pernah
memasang petunjuknya.

Karena itu, temuan yang berulang perlu diperiksa terhadap GDD, dan GDD itu
diperbarui. Dengan begitu, siklusnya benar-benar kembali ke awal, sebagaimana
diperlihatkan pada Gambar~\ref{fig:siklus-uji}.
